% Discussion
\section{Discussion}
\label{sec:discussion}

\subsection{Security Implication of the Floor Gap}

The result that matters for an operator is not how well either attack is detected at a rate the attacker would never choose, but how far each attacker can throttle down before detection stops being reliable. Section~\ref{sec:simulation} located that point directly: CP's 90\,\% crossing sits at roughly 44\,Hz, IFA's at roughly 1.8\,Hz, an order of magnitude apart, and Section~\ref{sec:mechanism} showed the gap is a property of what each attack is rather than of the network it runs over. Read as a bound on attacker behaviour rather than as a detector score, this says something asymmetric about the two threats. An IFA attacker can throttle down almost to the rate a single ordinary consumer already produces and still remain invisible, because a broken outcome is a broken outcome regardless of how rarely it occurs; there is close to no headroom between ``attacking'' and ``blending in.'' A CP attacker has no such option: to stay under the floor it must also stay under roughly the legitimate background's own volume, which caps how much cache damage it can do while quiet and forces it to choose between being loud and being effective. CP is therefore the easier attack to catch behaviorally, precisely because its only lever, added volume, is also its only liability. IFA is the more dangerous stealth threat of the two, not because it is harder to detect at any given rate, Table~\ref{tab:sim-floor} shows the opposite, but because the rate at which it stops being detectable is so close to the rate at which it stops being an attack at all. Crucially, this floor is measured against a realistic operating point rather than an idealized one: as Figure~\ref{fig:fpr} shows, the benign-consumer false-positive rate stays flat at a realistic 1--2\,\% across the swept attacker rates, rather than the trivial 0\,\% an idle network would give, so the floor reflects a genuine detection limit and not an artifact of a silent baseline.

\begin{figure}[t]
    \centering
    \includegraphics[width=0.7\textwidth]{fig_fpr_vs_rate}
    \caption{Benign-consumer false-positive rate against live traffic, swept across attacker rates. The rate stays a flat, realistic 1--2\,\%, rather than the trivial 0\,\% an idle network would produce.}
    \label{fig:fpr}
\end{figure}

\subsection{Comparison with Rule-Based Detection}

A conventional NDN security posture applies static thresholds to a handful of counters, typically interest rate, satisfaction ratio, PIT size, and NACK rate, and flags whichever counter crosses its line. Such a detector is structurally blind to Cache Pollution: CP's requests are answered, so satisfaction ratio, the metric most commonly monitored, stays high throughout the attack, and the only thing CP disturbs is a cache-hit ratio that a single-threshold rule rarely watches at all. The same detector is also weakened on its own home ground, IFA, by leaning on PIT size, which Section~\ref{sec:emulation} showed is non-discriminating once polling happens on a realistic multi-second interval rather than continuously. A rule-based detector's failure mode is therefore not narrow; it misses CP because it watches the wrong feature and misses IFA because the feature it does watch has already washed out by the time it is read. The Isolation Forest used throughout this paper does not need any single feature to cross a line, because it scores the joint improbability of the whole feature vector: a low cache-hit ratio co-occurring with otherwise ordinary traffic, or a collapsed satisfaction ratio alongside a normal-looking request rate, is exactly the kind of improbable combination a multi-feature model recognises and a single threshold cannot.

\subsection{Cross-Platform Consistency}

The floor result would carry less weight if it depended on one simulator's traffic model or one topology's geometry. It does not. On the emulation platform, Section~\ref{sec:method}'s global model, trained once across all three topologies, matches or beats per-topology models, and the same path-localised anomaly pattern, attacker and on-path nodes flagged, off-path nodes near baseline, is independently reproduced on both platforms: on miniNDN's real forwarding daemon (Section~\ref{sec:emulation}) and on ndnSIM's simulator (Section~\ref{sec:simulation}). Two platforms built on different codebases, one running actual NFD binaries under emulation and the other an ns-3 discrete-event model, reproducing the same qualitative behaviour is stronger evidence than either platform could offer alone; it suggests the pattern reflects how IFA propagates through NDN forwarding rather than an artifact of one simulator's approximations.

\subsection{A Note on Deployment}

To confirm the detector is usable outside a labelled offline evaluation, a real-time monitoring implementation, a live per-node scoring dashboard, was built and used to demonstrate the detector operating online; its architecture is outside this paper's scope.

\subsection{Limitations and Future Work}

Several limitations bound how these results should be read. The emulation track covers IFA only; CP required an active cache to leave a footprint at all (Section~\ref{sec:simulation}), so its evaluation rests entirely on simulation, and the two attacks were never run concurrently on either platform. miniNDN's 3\,s polling interval averages out sub-second PIT dynamics, which is the direct explanation for why \textit{pit\_size} and \textit{pit\_growth\_rate} proved non-discriminating in Section~\ref{sec:emulation} rather than evidence that PIT state is uninformative in principle; sub-second polling is a natural next step to test whether it recovers a usable PIT signal. The floor crossings reported in Section~\ref{sec:simulation} are interpolated between sampled rates rather than measured at the exact crossing point, so they are best read as bands, CP 40--45\,Hz, IFA 1--2\,Hz, and the two highest swept rates for each attack rest on a single seed rather than the five-seed average used in the floor region. The Isolation Forest deployed here is static: it does not adapt to gradual concept drift in legitimate traffic and would need periodic retraining in a long-running deployment. IFA's floor is set by traffic mimicry, the attacker leaves almost nothing outside the timeout and NACK ratios for the detector to work with, and augmenting the feature set with temporal statistics, for example autocorrelation of the NACK rate or of inter-arrival times, is a plausible route to pushing that floor lower still. Finally, none of this addresses content-integrity attacks that preserve normal traffic statistics while serving corrupted data; those attacks leave no behavioral trace for this kind of detector to find and would need complementary signature or provenance verification rather than anomaly detection.
